% One loop: a delay with its output fed back, three times over: the circuit you build, the loop it is, the sound you hear.
% The heavy stroke is the path that comes back: the feedback wire in the circuit, the circle on its own.
% Data: one-loop-harmonic-series.py (olhs-comb.dat). Everything here is computed (ink).
% Compile: lualatex one-loop-harmonic-series.tex ; SVG via PyMuPDF (dvisvgm's PDF route is broken on this Mac).
\documentclass[border=0pt]{standalone}
\usepackage{fontspec}
\setmainfont{JetBrains Mono}
\setmonofont{JetBrains Mono}
\usepackage{tikz}
\usepackage{pgfplots}
\usepackage{array}
\pgfplotsset{compat=1.16}
\usetikzlibrary{arrows.meta, calc, backgrounds, decorations.markings, positioning}
% Drafting Table, paper ground (_ops/drafting-table/design-system/colors_and_type.css)
\definecolor{paper}{HTML}{F0EEE5}
\definecolor{ink}{HTML}{1A1916}
\definecolor{ink2}{HTML}{3A3A35}
\definecolor{ink3}{HTML}{6A6862}
\definecolor{slab}{HTML}{D8D4C6}
\tikzset{
  >={Stealth[length=1.8mm]},
  loop/.style={draw=ink, line width=1.5pt},
  back/.style={draw=ink, line width=1.5pt, ->},
  wire/.style={draw=ink, line width=0.6pt},
  blk/.style={rectangle, draw=ink, fill=slab, line width=0.7pt, minimum height=7mm, minimum width=12mm, inner sep=2pt, font=\scriptsize, text=ink},
  adder/.style={circle, draw=ink, fill=paper, line width=0.7pt, minimum size=4.6mm, inner sep=0pt, font=\scriptsize},
  tap/.style={circle, fill=ink, minimum size=3pt, inner sep=0pt},
  lab/.style={font=\scriptsize, text=ink2},
  rowhead/.style={font=\bfseries\small, anchor=west, text=ink},
  rowsub/.style={font=\scriptsize, anchor=west, text=ink2},
  colhead/.style={font=\bfseries\scriptsize, anchor=west, text=ink},
  spin/.style={postaction={decorate}, decoration={markings, mark=at position 0.25 with {\arrow{>}}}},
}
\pgfplotsset{
  /pgf/number format/assume math mode=true,
  snd/.style={width=6.6cm, height=3.0cm, axis lines=left, tick align=outside, ymin=0, ytick=\empty,
    axis line style={ink, line width=0.5pt}, tick style={ink}, clip=false,
    xlabel style={font=\scriptsize, text=ink2, yshift=1mm}, ylabel style={font=\scriptsize, text=ink2},
    tick label style={font=\tiny, text=ink2}},
}
% a row: frame (top y, height), heading, subtitle, and the two column rules
\newcommand{\rowframe}[4]{%
  \draw[ink, line width=0.5pt] (-0.3,#1) rectangle (19.9,#1-#2);
  \draw[ink, line width=0.3pt] (6.9,#1-0.95) -- (6.9,#1-#2) (12.4,#1-0.95) -- (12.4,#1-#2) (-0.3,#1-0.95) -- (19.9,#1-0.95);
  \node[rowhead] at (0,#1-0.32) {#3};
  \node[rowsub] at (0,#1-0.7) {#4};}
\begin{document}
\begin{tikzpicture}

% ---------- column heads ----------
\node[colhead] at (0,1.15) {THE CIRCUIT \quad what you build};
\node[colhead] at (7.2,1.15) {THE LOOP \quad what it is};
\node[colhead] at (12.7,1.15) {THE SOUND \quad what you hear};
\node[font=\scriptsize, anchor=west, text=ink2] at (0,0.55) {Heavy: the path that comes back. In the circuit it is the feedback wire; on its own it is a circle.};

% =============== ONE LOOP ===============
\def\y{0}
\rowframe{\y}{7.3}{ONE LOOP}{one loop, one fundamental and its whole-number multiples: every k / T}
\begin{scope}[shift={(0,\y)}]
  % the circuit
  \node[lab] (in) at (0.1,-3.0) {in};
  \node[adder] (a) at (0.9,-3.0) {+};
  \node[blk] (d) at (2.6,-3.0) {delay T};
  \node[tap] (t) at (4.3,-3.0) {};
  \draw[wire,->] (in) -- (a); \draw[wire,->] (a) -- (d); \draw[wire] (d) -- (t);
  \draw[wire,->] (t) -- (6.65,-3.0) node[lab, above left] {out};
  \node[blk, minimum width=9mm] (g) at (2.6,-4.6) {× g};
  \draw[back] (t) |- (g.east);
  \draw[back] (g.west) -| (a);
  \node[lab, anchor=north, align=center] at (3.2,-5.4) {what leaves the delay\\comes back to the start};

  % the loop: one circle, one trip of length T
  \draw[loop, spin] (8.3,-3.0) circle (1.0); \node[font=\small, text=ink] at (8.3,-3.0) {T};
  \fill[ink] ($(8.3,-3.0)+(-40:1.0)$) circle (2.2pt);
  \node[lab, anchor=north, align=center] at (8.3,-4.35) {one trip round\\takes T};
  % a wavetable phasor is the same loop, unrolled
  \begin{scope}[shift={(10.0,-3.55)}]
    \draw[wire, ink3, ->] (0,0) -- (1.85,0);
    \draw[wire, ink3, ->] (0,0) -- (0,1.3);
    \node[font=\tiny, text=ink2, anchor=east] at (0,1.0) {1};
    \node[font=\tiny, text=ink2, anchor=east] at (0,0) {0};
    \draw[loop] (0,0) -- (0.75,1.0); \draw[ink3, line width=0.4pt, densely dashed] (0.75,1.0) -- (0.75,0);
    \draw[loop] (0.75,0) -- (1.5,1.0); \draw[ink3, line width=0.4pt, densely dashed] (1.5,1.0) -- (1.5,0);
    \draw[ink3, line width=0.3pt] (0,-0.15) -- (0,-0.25) -- (0.75,-0.25) -- (0.75,-0.15);
    \node[font=\tiny, text=ink2, anchor=north] at (0.375,-0.25) {T};
    \fill[ink] (0.75,1.0) circle (1.6pt);
  \end{scope}
  \node[lab, anchor=north, align=center] at (10.95,-4.35) {a wavetable\\phasor is the\\same loop};

  % the sound
  \begin{axis}[snd, at={(12.85cm,-4.15cm)}, height=3.2cm, xmin=0, xmax=6.5, ymax=5.6, xtick={0,...,6},
      xlabel={frequency × T \quad (peaks at k)}]
    \addplot[line width=1.0pt, ink] table[x=fT, y=mag] {olhs-comb.dat};
  \end{axis}
  \node[anchor=north west, inner sep=0pt] at (12.75,-5.2) {%
    \begin{tabular}{@{}>{\centering\arraybackslash}p{2.35cm}@{}>{\centering\arraybackslash}p{0.4cm}@{}>{\centering\arraybackslash}p{1.8cm}@{}>{\centering\arraybackslash}p{0.4cm}@{}>{\centering\arraybackslash}p{2.0cm}@{}}
      \small\bfseries f & \small\bfseries = & \small\bfseries k & \small\bfseries / & \small\bfseries T \\[3pt]
      \scriptsize\color{ink2}(a frequency the loop can hold) & \scriptsize\color{ink2}= & \scriptsize\color{ink2}(any whole number) & \scriptsize\color{ink2}÷ & \scriptsize\color{ink2}(the time of one trip) \\
    \end{tabular}};
\end{scope}
\node[font=\scriptsize, anchor=north west, text=ink2, text width=19.6cm] at (0,-7.55) {A cycle the loop can hold arrives back in step with itself after one trip, so its period divides T: T, T/2, T/3 \ldots{} Each trip adds it again and it grows; everything else meets itself out of step and fades. Computed with g = 0.8: peaks 1/(1 − g) = 5, dips 1/(1 + g) ≈ 0.56.};

% ---------- title block ----------
\draw[ink, line width=0.8pt] (-0.3,-8.65) rectangle (19.9,-9.55);
\draw[ink, line width=0.3pt] (9.6,-8.65) -- (9.6,-9.55) (14.6,-8.65) -- (14.6,-9.55);
\node[font=\scriptsize\bfseries, text=ink, anchor=west] at (-0.1,-9.1) {DSP IN LOOPING DIMENSIONS · ONE LOOP};
\node[font=\scriptsize, text=ink, anchor=west] at (9.8,-9.1) {INK computed};
\node[font=\scriptsize, text=ink2, anchor=west] at (14.8,-9.1) {FIG · 2026-11-09 · PAPER};

\begin{scope}[on background layer]
  \fill[paper] ($(current bounding box.south west)+(-0.4,-0.4)$) rectangle ($(current bounding box.north east)+(0.4,0.4)$);
\end{scope}
\end{tikzpicture}
\end{document}
